%% Blank Elsevier (elsarticle) article skeleton, trimmed to fields this app collects.
%%
%% Base mode is "preprint" (review-stage, single column, wide spacing) -
%% matches this app's submission-stage focus. "final" mode (camera-ready)
%% requires also picking a page-format preset (1p/3p/5p - see below);
%% not something this app currently asks the user to choose.
%%
%% twocolumn is a separate toggle from the page-format preset, and maps
%% to FormData.columnLayout (add "twocolumn" for double, omit for single).
%% authoryear/num class options map to FormData.referencingStyle
%% (authoryear for APA/Harvard, num for IEEE/Vancouver-style numbering).
\documentclass[preprint,12pt,authoryear]{elsarticle}

\usepackage{amssymb}
\usepackage{amsmath}

%% Journal name - not currently a FormData field; leave as a placeholder
%% or drop this line entirely.
\journal{}

\begin{document}

\begin{frontmatter}

%% Title -> FormData.title
\title{}

%% One \author{}\ead{}\affiliation{} block per entry in FormData.authors -
%% repeat for each author, do not hardcode a fixed number.
%% \ead{} is the email command (FormData.authors[i].email).
%% For the author where corresponding === true, add \corref{cor1} after
%% \author{Name} and a matching \cortext[cor1]{Corresponding author}.
%% elsarticle has no native ORCID field; if FormData.authors[i].orcidId
%% is provided, append it via \fnref{}/\fntext{}, e.g. "ORCID: 0000-...".
\author{}
\ead{}

\affiliation{organization={},
            addressline={},
            city={},
            postcode={},
            state={},
            country={}}

%% Abstract -> FormData.abstract
\begin{abstract}

\end{abstract}

%% Highlights - Elsevier-specific; only include this block when
%% FormData.highlights === "yes" (this app already gates the Highlights
%% field to Elsevier only). No dedicated FormData field holds the
%% highlight text itself - use short bullets summarizing FormData.abstract
%% or FormData.conclusion.
\begin{highlights}
\item
\item
\end{highlights}

%% Keywords -> FormData.keywords. Note: elsarticle requires \sep between
%% keyword entries in the .tex source itself, regardless of what
%% FormData.keywordSeparator was set to for display purposes elsewhere.
\begin{keyword}

\end{keyword}

\end{frontmatter}

%% One \section{} per entry the user checked in FormData.keySections
%% (Introduction, Literature Review, Methodology, Experimentation,
%% Results, Discussion). Generate only the sections that were checked -
%% do not invent additional ones. Always include a Conclusion section
%% driven by FormData.conclusion.
\section{Introduction}
\label{sec:intro}

%% Placeholder table (always include exactly one, in this section only) -
%% lets caption/table font size fields be visually verified in the
%% compiled output.
\begin{table}[t]
\centering
\begin{tabular}{l c r}
  Column A & Column B & Column C \\
  -- & -- & -- \\
  -- & -- & -- \\
\end{tabular}
\caption{Placeholder table caption.}\label{tab:placeholder}
\end{table}

\section{Conclusion}
\label{sec:conclusion}

%% Appendix - not tied to a FormData field. elsarticle uses the bare
%% \appendix command; sections after it are lettered, not numbered.
\appendix
\section{}
\label{app:1}

\subsection{}
\label{app:1.1}

%% Bibliography style -> FormData.referencingStyle (elsarticle-harv for
%% author-year, elsarticle-num for numbered). If FormData.bibliographyText
%% or an uploaded .bib file (FormData.bibliography) is present, keep
%% \bibliographystyle{}+\bibliography{}; otherwise use a manual
%% thebibliography block with a placeholder entry matching the chosen style.
\begin{thebibliography}{00}

\bibitem{}

\end{thebibliography}

\end{document}
%% End of trimmed elsarticle skeleton.